\documentclass[10pt,a4paper]{amsart}
\usepackage[margin=19mm]{geometry}
\usepackage{amsmath,amssymb,mathtools}
\usepackage[scr=rsfs,cal=euler]{mathalfa}
\usepackage{xfrac}
\usepackage[unicode,hidelinks,bookmarks=false]{hyperref}
\newcommand{\ie}{i.e.\ }
\newcommand{\cf}{cf.\ }
\newcommand{\resp}{respectively\ }
\newcommand{\Subsection}{\subsection}
\providecommand{\nobreakdash}{\nobreak\hbox{-}}
\setlength{\emergencystretch}{2em}
\multlinegap0pt
\usepackage{amssymb}
\usepackage{algorithm}\usepackage{algpseudocode}
\usepackage{tikz}
\newtheorem{lemma}{Lemma}
\title{Contraction Analysis of Holomorphic Dynamical Systems via the Intrinsic Kobayashi Metric}
\author{Soumic Sarkar}
\date{Source: arXiv:2608.07551v1 (CC BY 4.0)}
\begin{document}
\maketitle

\noindent\emph{Source and licence.} Contraction Analysis of Holomorphic Dynamical Systems via the Intrinsic Kobayashi Metric. Version arXiv:2608.07551v1 (CC BY 4.0), distributed by arXiv under the Creative Commons Attribution 4.0 International licence, http://creativecommons.org/licenses/by/4.0/, as recorded in the arXiv licence metadata for this exact version and shown on the abstract page https://arxiv.org/abs/2608.07551v1. Full licence notice: \texttt{licenses/arxiv-2608.07551-CC-BY-4.0.txt}.\par\smallskip

\noindent\textit{Editorial scope.} Counted unit: the article's lemma lem:dini (source0.tex lines 647--650). The article's complete original proof of it is delivered with it (source0.tex lines 652--672, one \texttt{proof} environment of 21 lines). No mathematics is written, repaired or re-derived: the statement and the proof are the article's own lines, in the article's own notation, and no retained line is substituted or re-flowed.  No further source-local context is needed: every cross-reference of every retained line resolves inside the two delivered spans.  Nothing was omitted from any retained span, so the package carries no omission record.  No retained line contains a citation, so the wrapper carries no bibliography and no bibliography entry.  No retained line contains a figure environment, an image-inclusion command or drawing code, so no image is delivered and the package declares no asset dependency.  Editorial formatting only, in addition: an AMS \texttt{amsart} preamble that restates the article's own declarations and macros used by the retained lines, namely the environments \texttt{lemma} on separate counters, together with the standard packages those lines need.  The retained environments print in this document as Lemma 1 in order of appearance, whereas the article prints the same items with the numbering of its own full document.  The complete native source of the article as distributed in the arXiv e-print for this exact version is bundled unchanged as \texttt{sources/207091/source0.tex}.
\begin{lemma}
\label{lem:dini}
Let $\psi:[0,\infty)\to\mathbb{R}_{\ge0}$ be continuous and suppose $D^+\psi(t)\le-\lambda\psi(t)$ for all $t\ge0$ and some $\lambda>0$. Then $\psi(t)\le e^{-\lambda t}\psi(0)$ for all $t\ge0$.
\end{lemma}

\begin{proof}
Let $\varphi(t)=e^{\lambda t}\psi(t)$; $\varphi$ is continuous since $\psi$ is. For $\epsilon>0$,
\begin{equation}
\label{eq:dinistep}
\begin{aligned}
\frac{\varphi(t+\epsilon)-\varphi(t)}{\epsilon}
&=
e^{\lambda t}\Bigg(
e^{\lambda\epsilon}
\frac{\psi(t+\epsilon)-\psi(t)}{\epsilon}
\\
&\qquad\qquad
+\psi(t+\epsilon)
\frac{e^{\lambda\epsilon}-1}{\epsilon}
\Bigg).
\end{aligned}
\end{equation}
Taking $\limsup_{\epsilon\to0^+}$ of \eqref{eq:dinistep} and using continuity of $\psi$ at $t$ to evaluate the second term in the limit gives $D^+\varphi(t)=e^{\lambda t}\big(D^+\psi(t)+\lambda\psi(t)\big)\le0$ for every $t\ge0$, using the hypothesis on $\psi$.

It remains to show that a continuous function with nonpositive upper Dini derivative everywhere on an interval is non-increasing on that interval. Suppose not: then $\varphi(t_2)>\varphi(t_1)$ for some $0\le t_1<t_2$. Let $c=\varphi(t_1)$ and let $\tau=\sup\{t\in[t_1,t_2]:\varphi(t)\le c\}$. By continuity $\varphi(\tau)=c$ and $\tau<t_2$, and by definition of $\tau$ as a supremum, $\varphi(s)>c$ for $s$ in some sequence approaching $\tau$ from above, so $\limsup_{\epsilon\to0^+}(\varphi(\tau+\epsilon)-\varphi(\tau))/\epsilon\ge0$; in fact $\varphi(s)>c=\varphi(\tau)$ for all $s\in(\tau,t_2]$ by definition of $\tau$, so this quotient is strictly positive along that sequence and $D^+\varphi(\tau)>0$, contradicting $D^+\varphi\le0$ everywhere. Hence $\varphi$ is non-increasing, so $\varphi(t)\le\varphi(0)$, which is exactly $\psi(t)\le e^{-\lambda t}\psi(0)$.
\end{proof}

\end{document}
